\slajdid{09_1-s050}
\begin{frame}[label=09_1-s050]
\gusttekst\begin{center}\input{slike/tikz/s050_invertor_gejt_na_napajanju.tex}\end{center}
Za ulazni napon $V_I<V_{Tn,D}$ tranzistor $T_D$ je zakočen. Kolo je neoptrećeno $I_{Dn,D}=0=I_{Dn,L}$. Kako tranzistor $T_L$ radi u zasićenju
\[I_{Dn,L}=\frac{k_{n,L}}2(V_{GS,L}-V_{Tn,L})^2=0\]
njegova radna tačka će se podesiti tako da je $V_{GS,L}=V_{Tn,L}$. To je jedina moguća radna tačka, Prema tome
\[V_O=V_{OH}=V_{DD}-V_{GS,L}=V_{DD}-V_{Tn,L}\]
U odnosu na prethodni slučaj napon logičke jedinice je manji
\end{frame}
